\section{Determinization}
We use a type-directed transformation that preserves non-probabilistic structure and replaces probabilistic sampling by its expected value when the result is typed at expectation mode. \Cref{fig:determinization} shows the determinization function, defined inductively on the structure of $e : \tau$. For all other language constructs, $\exptrans{\cdot}$ recurses structurally over them.

\begin{figure}[t]
  \centering
  \tiny
  \setlength{\tabcolsep}{6pt}
  \renewcommand{\arraystretch}{1.2}
  \begin{adjustbox}{max width=\textwidth}
  \begin{tabular}{@{}>{\raggedright\arraybackslash}m{0.33\textwidth}
      @{\hspace{3em}}
      >{\raggedright\arraybackslash}p{0.14\textwidth}
      @{\hspace{1em}}
      >{\raggedright\arraybackslash}p{0.43\textwidth}@{}}
    \toprule
    \toprule
    \textbf{$e:\tau$} & \textbf{Condition} & \textbf{$\exptrans{e : \tau}$} \\
    \midrule
    %
    $x : \tau$  & 
    -- &
    $x$ \\
    \midrule
    %
    $c : \Float{m}$ &
    -- &
    $c$ \\
    \midrule
    %
    $\uniform(e_1 : \Float{m_1}, e_2 : \Float{m_2}) : \Float{m_3}$ &
    \begin{tabular}[c]{@{}l@{}}
      $m_3 = \E$ \\
      $m_3 = \G$
    \end{tabular} &
    \begin{tabular}[c]{@{}l@{}}
      $(\exptrans{e_1 : \Float{m_1}} + \exptrans{e_2 : \Float{m_2}}) \times 0.5$ \\
      $\uniform(\exptrans{e_1 : \Float{m_1}}, \exptrans{e_2 : \Float{m_2}})$
    \end{tabular} \\
    \midrule
    %
    %
    $\gaussian(e_1 : \Float{m_1}, e_2 : \Float{m_2}) : \Float{m_3}$ &
    \begin{tabular}[c]{@{}l@{}}
      $m_3 = \E$ \\
      $m_3 = \G$
    \end{tabular} &
    \begin{tabular}[c]{@{}l@{}}
      $\exptrans{e_1 : \Float{m_1}}$ \\
      $\gaussian(\exptrans{e_1 : \Float{m_1}}, \exptrans{e_2 : \Float{m_2}})$
    \end{tabular} \\
    \midrule
    %
    $\betafn(e_1 : \Float{m_1}, e_2 : \Float{m_2}) : \Float{m_3}$ &
    \begin{tabular}[c]{@{}l@{}}
      $m_3 = \E$ \\
      $m_3 = \G$
    \end{tabular} &
    \begin{tabular}[c]{@{}l@{}}
      $\exptrans{e_1 : \Float{m_1}} / (\exptrans{e_1 : \Float{m_1}} + \exptrans{e_2 : \Float{m_2}})$ \\
      $\betafn(\exptrans{e_1 : \Float{m_1}}, \exptrans{e_2 : \Float{m_2}})$
    \end{tabular} \\
    \midrule
    %
    $\poisson(e_1 : \Float{m_1}) : \Float{m_2}$ &
    \begin{tabular}[c]{@{}l@{}}
      $m_2 = \E$ \\
      $m_2 = \G$
    \end{tabular} &
    \begin{tabular}[c]{@{}l@{}}
      $\exptrans{e_1 : \Float{m_1}}$ \\
      $\poisson(\exptrans{e_1 : \Float{m_1}})$
    \end{tabular} \\
    \midrule
    %
    $\gammafn(e_1 : \Float{m_1}, e_2 : \Float{m_2}) : \Float{m_3}$ &
    \begin{tabular}[c]{@{}l@{}}
      $m_3 = \E$ \\
      $m_3 = \G$
    \end{tabular} &
    \begin{tabular}[c]{@{}l@{}}
      $\exptrans{e_1 : \Float{m_1}} / \exptrans{e_2 : \Float{m_2}}$ \\
      $\gammafn(\exptrans{e_1 : \Float{m_1}}, \exptrans{e_2 : \Float{m_2}})$
    \end{tabular} \\
    \midrule
    %
    $\exponential(e_1 : \Float{m_1}) : \Float{m_2}$ &
    \begin{tabular}[c]{@{}l@{}}
      $m_2 = \E$ \\
      $m_2 = \G$
    \end{tabular} &
    \begin{tabular}[c]{@{}l@{}}
      $1 / \exptrans{e_1 : \Float{m_1}}$ \\
      $\exponential(\exptrans{e_1 : \Float{m_1}})$
    \end{tabular} \\
    \midrule
    %
    $(e_1 : \Float{m_1} + e_2 : \Float{m_2}) : \Float{m_3}$ &
    -- &
    $\exptrans{e_1 : \Float{m_1}} + \exptrans{e_2 : \Float{m_2}}$ \\
    \midrule
    %
    $(e_1 : \Float{m_1} \times e_2 : \Float{m_2}) : \Float{m_3}$ &
    -- &
    $\exptrans{e_1 : \Float{m_1}} \times \exptrans{e_2 : \Float{m_2}}$ \\
    \midrule
    %
    $(\ifkw\ b : \Bool\ \thenkw\ e_1 : \tau\ \elsekw\ e_2 : \tau) : \tau$ &
    -- &
    $\ifkw\ \exptrans{b : \Bool}\ \thenkw\ \exptrans{e_1 : \tau}\ \elsekw\ \exptrans{e_2 : \tau}$ \\
    %
    \midrule
    $(\letkw\ x = e_1 : \tau_1\ \inkw\ e_2 : \tau_2) : \tau_2$ &
    -- &
    $\letkw\ x = \exptrans{e_1 : \tau_1}\ \inkw\ \exptrans{e_2 : \tau_2}$
    \\
    \bottomrule
    \bottomrule
  \end{tabular}
  \end{adjustbox}
  \caption{Determinization function. We case split only on the modes that affect $\exptrans{\cdot}$; all other float modes are independent and may range over $\E$ and $\G$.}
  \label{fig:determinization}
\end{figure}

\begin{example}
  Recall the following program, annotated with types:
  \[
    \begin{aligned}
      &(\letkw\ x = \uniform(0 : \Float{\G},\, 1 : \Float{\G}) : \Float{\G} \\
      &\quad \inkw\ \betafn(x : \Float{\G},\, 2 : \Float{\G}) : \Float{\E}) : \Float{\E}.
    \end{aligned}
  \]

  Based on determinization, obtain:
  \[
    \letkw\ x = \uniform(0,1)\ \inkw\ x / (x + 2).
  \]
\end{example}

\begin{example}
  Recall the following program, annotated with types:
  \[
    \begin{aligned}
    &(\ifkw\ \flip(0.5) \\
    &\quad \thenkw\ \gammafn\!\big(
      \uniform(0 : \Float{\E},\, 1 : \Float{\E}) : \Float{\E},\,
      1 : \Float{\G}
    \big) : \Float{\E} \\
    &\quad \elsekw\ \big(
      2 : \Float{\E}
      *\,
      \betafn(1 : \Float{\G},\, 4 : \Float{\G}) : \Float{\E}
    \big) : \Float{\E} \\
    &)\ : \Float{\E}.
    \end{aligned}
  \]

  Based on determinization, obtain:
  \[
    \ifkw\ \flip( 0.5 )\ \thenkw\ (( 0 + 1 ) * 0.5) / 1\ \elsekw\ 2 * ( 1 / ( 1 + 4 ) )
  \]
\end{example}


% We use a type-directed transformation \(\Gamma \vdash e : \tau \darrow e'\) that preserves non-probabilistic structure and replaces probabilistic sampling by its expected value when the result is typed at expectation mode. Representative clauses are below; other constructs follow the obvious structural recursion.
% \begin{mathpar}
%   \inferrule[\textsc{Det-Uniform-E}]
%     {\Gamma \vdash e_1 : \Float{\E} \darrow e_1' \\ \Gamma \vdash e_2 : \Float{\E} \darrow e_2'}
%     {\Gamma \vdash \uniform(e_1,e_2) : \Float{\E} \darrow (e_1' + e_2') \times 0.5}
%   \and
%   \inferrule[\textsc{Det-Uniform-G}]
%     {\Gamma \vdash e_1 : \Float{\G} \darrow e_1' \\ \Gamma \vdash e_2 : \Float{\G} \darrow e_2'}
%     {\Gamma \vdash \uniform(e_1,e_2) : \Float{\G} \darrow \uniform(e_1',e_2')}
%   \and
%   \inferrule[\textsc{Det-Gauss-E}]
%     {\Gamma \vdash e_1 : \Float{\E} \darrow e_1' \\ \Gamma \vdash e_2 : \Float{\G} \darrow e_2'}
%     {\Gamma \vdash \gaussian(e_1,e_2) : \Float{\E} \darrow e_1'}
%   \and
%   \inferrule[\textsc{Det-Gauss-G}]
%     {\Gamma \vdash e_1 : \Float{\G} \darrow e_1' \\ \Gamma \vdash e_2 : \Float{\G} \darrow e_2'}
%     {\Gamma \vdash \gaussian(e_1,e_2) : \Float{\G} \darrow \gaussian(e_1',e_2')}
%   \and
%   \inferrule[\textsc{Det-Var}]
%     {\Gamma \vdash x : \tau}
%     {\Gamma \vdash x : \tau \darrow x}
%   \and
%   \inferrule[\textsc{Det-Abs}]
%     {\Gamma, x{:}\tau_1 \vdash e : \tau_2 \darrow e'}
%     {\Gamma \vdash \lambda x.\,e : \tau_1 \to \tau_2 \darrow \lambda x.\,e'}
%   \and
%   \inferrule[\textsc{Det-Rec}]
%     {\Gamma, f{:}\tau_1 \to \tau_2, x{:}\tau_1 \vdash e : \tau_2 \darrow e'}
%     {\Gamma \vdash \mathsf{rec}\;f\;x.\,e : \tau_1 \to \tau_2 \darrow \mathsf{rec}\;f\;x.\,e'}
%   \and
%   \inferrule[\textsc{Det-Const}]
%     {\Gamma \vdash c : \Float{m}}
%     {\Gamma \vdash c : \Float{m} \darrow c}
%   \and
%   \inferrule[\textsc{Det-Neg}]
%     {\Gamma \vdash e : \Float{m} \darrow e'}
%     {\Gamma \vdash -e : \Float{m} \darrow -e'}
%   \and
%   \inferrule[\textsc{Det-Add}]
%     {\Gamma \vdash e_1 : \Float{m} \darrow e_1' \\ \Gamma \vdash e_2 : \Float{m} \darrow e_2'}
%     {\Gamma \vdash e_1 + e_2 : \Float{m} \darrow e_1' + e_2'}
%   \and
%   \inferrule[\textsc{Det-Mul}]
%     {\Gamma \vdash e_1 : \Float{\G} \darrow e_1' \\ \Gamma \vdash e_2 : \Float{\G} \darrow e_2'}
%     {\Gamma \vdash e_1 \times e_2 : \Float{\G} \darrow e_1' \times e_2'}
%   \and
%   \inferrule[\textsc{Det-If}]
%     {\Gamma \vdash e : \Bool \darrow e' \\ \Gamma \vdash e_1 : \tau \darrow e_1' \\ \Gamma \vdash e_2 : \tau \darrow e_2'}
%     {\Gamma \vdash \ifkw\;e\;\thenkw\;e_1\;\elsekw\;e_2 : \tau \darrow \ifkw\;e'\;\thenkw\;e_1'\;\elsekw\;e_2'}
% \end{mathpar}
% Lambda, application, pairs, sums, and \(\letkw\) follow the same shape (recursing on subterms). The key point is that expectation-mode samples are determinized to their expected value, while general-mode samples remain probabilistic but have their subterms determinized.

